\chapter{Connexion et Profil Utilisateur}
\label{chap:connexion}

\begin{rolebox}{Tous les utilisateurs}
  Ce chapitre s'applique \`a l'ensemble des profils. Chaque utilisateur doit se connecter avant d'acc\'eder \`a l'application.
\end{rolebox}

\section{Page de connexion}

\`A l'ouverture de l'application, l'utilisateur arrive sur la page de connexion. Aucun acc\`es aux donn\'ees n'est possible sans authentification.

\screenshot{img_login}{Page de connexion}{Capturer la page de connexion de l'application dans son int\'egralit\'e. Le formulaire avec le champ identifiant et le bouton de validation doit \^etre clairement visible. Fen\^etre en plein \'ecran, fond de page visible.}

\section{Proc\'edure de connexion}

\begin{steps}
  \item Ouvrir le navigateur web et saisir l'adresse de l'application fournie par l'administrateur.
  \item Sur la page de connexion, saisir votre \textbf{identifiant} (login) dans le champ pr\'evu.
  \item Cliquer sur le bouton \textbf{Se connecter} (ou appuyer sur la touche \textsf{Entr\'ee}).
  \item L'application v\'erifie votre identit\'e et charge votre profil, vos affectations et l'ann\'ee acad\'emique courante.
  \item Vous \^etes redirig\'e automatiquement vers le \textbf{Tableau de bord}.
\end{steps}

\begin{notebox}[Syst\`eme d'authentification]
  L'application utilise un syst\`eme de connexion par \textbf{identifiant unique} (login). Selon la configuration de votre \'etablissement, cet identifiant peut correspondre \`a votre identifiant institutionnel (num\'ero de personnel, alias de messagerie, etc.). Rapprochez-vous de votre administrateur en cas de probl\`eme de connexion.
\end{notebox}

\begin{warnbox}[Identifiant incorrect]
  Si votre identifiant n'est pas reconnu, l'application affiche un message d'erreur. V\'erifiez que\,: (1) vous avez saisi l'identifiant correct, (2) votre compte a bien \'et\'e cr\'e\'e par un administrateur, (3) votre compte poss\`ede au moins une affectation active dans l'ann\'ee EN\_COURS.
\end{warnbox}

\section{Tableau de bord apr\`es connexion}

\`A la connexion, l'application affiche le tableau de bord. Celui-ci pr\'esente\,:

\begin{itemize}
  \item \textbf{L'ann\'ee acad\'emique active} visible dans l'en-t\^ete (ex. : \og 2024-2025 \fg).
  \item Les \textbf{raccourcis} vers les principales fonctionnalit\'es accessibles selon votre r\^ole.
  \item Un \textbf{indicateur de notifications} (cloche) signalant les \'ev\`enements en attente.
  \item Pour les \textbf{Services Centraux uniquement} : un s\'electeur d'ann\'ee permettant de changer de contexte.
\end{itemize}

\screenshot{img_dashboard_general}{Tableau de bord apr\`es connexion}{Capturer la page du tableau de bord telle qu'elle appara\^it apr\`es connexion avec un compte de test. Montrer l'en-t\^ete (ann\'ee courante, cloche de notifications), le menu de navigation lat\'eral et les cartes/tuiles du tableau de bord. Plein \'ecran.}

\section{S\'electeur d'ann\'ee acad\'emique}

\begin{rolebox}{services-centraux uniquement}
  Le changement d'ann\'ee est r\'eserv\'e au r\^ole \texttt{services-centraux}.
\end{rolebox}

Pour les utilisateurs ayant le r\^ole \texttt{services-centraux}, un menu d\'eroulant dans l'en-t\^ete permet de s\'electionner n'importe quelle ann\'ee (EN\_COURS, PREPARATION, ARCHIVEE).

\begin{steps}
  \item Localiser le s\'electeur d'ann\'ee dans l'en-t\^ete de l'application (g\'en\'eralement en haut \`a droite).
  \item Cliquer sur l'ann\'ee actuellement affich\'ee pour d\'eplier la liste des ann\'ees disponibles.
  \item Cliquer sur l'ann\'ee souhait\'ee.
  \item L'application recharge automatiquement toutes les donn\'ees dans le contexte de l'ann\'ee s\'electionn\'ee.
\end{steps}

\screenshot{img_selecteur_annee}{S\'electeur d'ann\'ee (Services Centraux)}{Capturer l'en-t\^ete de l'application avec le menu d\'eroulant du s\'electeur d'ann\'ee ouvert, montrant plusieurs ann\'ees disponibles. Utiliser un compte services-centraux.}

\begin{notebox}
  Pour tous les autres r\^oles, l'ann\'ee affich\'ee est l'ann\'ee \textbf{EN\_COURS} et ne peut pas \^etre modifi\'ee. Elle est visible mais non interactive.
\end{notebox}

\section{Navigation dans l'application}

L'application dispose d'un menu de navigation (lat\'eral ou en haut de page) dont le contenu varie selon le r\^ole de l'utilisateur.

Les rubriques possibles sont\,:

\medskip
\begin{tabularx}{\textwidth}{L{4cm}L{5cm}X}
  \toprule
  \rowcolor{cgplightblue}
  \textbf{Rubrique} & \textbf{Ic\^one / Label} & \textbf{Accessible par} \\
  \midrule
  Tableau de bord   & \faHome\ Tableau de bord  & Tous \\
  \rowcolor{cgplightgray}
  Rechercher        & \faSearch\ Rechercher      & Tous \\
  Organigramme      & \faProjectDiagram\ Organigramme & Tous \\
  \rowcolor{cgplightgray}
  Import / Export   & \faFileExcel\ Import/Export & SC, Directeurs composante/admin \\
  Responsables      & \faUsers\ Responsables    & SC, Directeurs, Admins \\
  \rowcolor{cgplightgray}
  Structures        & \faBuilding\ Structures   & SC, Directeurs composante \\
  D\'emandes de r\^oles & \faClipboardList\ D\'emandes & SC, Directions \\
  \rowcolor{cgplightgray}
  D\'el\'egations   & \faUserCheck\ D\'el\'egations & SC, Directeurs \\
  Ann\'ees          & \faCalendar\ Ann\'ees     & SC uniquement \\
  \rowcolor{cgplightgray}
  Audit             & \faHistory\ Audit         & SC uniquement \\
  Signalements      & \faFlag\ Signalements     & Tous \\
  \rowcolor{cgplightgray}
  Mon profil        & \faUser\ Profil           & Tous \\
  \bottomrule
\end{tabularx}

\medskip
\begin{notebox}
  Les rubriques non accessibles pour un r\^ole donn\'e ne s'affichent pas dans le menu. Le menu est donc adapt\'e automatiquement.
\end{notebox}

\screenshot{img_menu_navigation}{Menu de navigation (exemple Services Centraux)}{Capturer le menu de navigation complet avec toutes les rubriques visibles pour un compte services-centraux (c'est le menu le plus complet). Montrer le menu lat\'eral enti\`erement d\'eploy\'e.}

\section{D\'econnexion}

Pour quitter l'application, cliquer sur l'ic\^one de profil ou le menu utilisateur (g\'en\'eralement en haut \`a droite de l'\'ecran), puis s\'electionner \textbf{D\'econnecter} ou \textbf{Se d\'econnecter}.

\begin{warnbox}[Session navigateur]
  L'application conserve votre identifiant de session dans le navigateur (\texttt{localStorage}). Si vous partagez votre poste de travail, pensez \`a vous d\'econnecter explicitement apr\`es chaque utilisation.
\end{warnbox}

\section{Gestion du profil}

La rubrique \textbf{Mon Profil} (accessible depuis le menu de navigation ou l'ic\^one utilisateur) vous permet de consulter et de modifier une partie de vos informations personnelles.

\subsection{Consultation du profil}

La fiche profil affiche\,:
\begin{itemize}
  \item \textbf{Informations d'identit\'e} : nom, pr\'enom (non modifiables par l'utilisateur).
  \item \textbf{Identifiant de connexion} (non modifiable).
  \item \textbf{R\^oles et affectations} : liste des r\^oles actifs pour l'ann\'ee courante.
  \item \textbf{Coordonn\'ees modifiables} : courriel, t\'el\'ephone, bureau.
\end{itemize}

\screenshot{img_profil}{Page Mon Profil}{Capturer la page profil utilisateur enti\`ere. Montrer les sections informations d'identit\'e, les affectations actives et les champs modifiables (courriel, t\'el\'ephone, bureau). Utiliser des donn\'ees de test non confidentielles.}

\subsection{Modification du profil}

\begin{steps}
  \item Naviguer vers \textbf{Mon Profil} dans le menu.
  \item Dans la section \og Mes coordonn\'ees \fg, cliquer sur \textbf{Modifier}.
  \item Mettre \`a jour les champs souhait\'es\,: \textit{Courriel secondaire}, \textit{T\'el\'ephone}, \textit{Bureau}.
  \item Cliquer sur \textbf{Enregistrer}.
\end{steps}

\begin{warnbox}[Champs non modifiables]
  Certains champs ne peuvent \^etre modifi\'es que par un administrateur ou les Services Centraux\,: \textit{Nom}, \textit{Pr\'enom}, \textit{Identifiant}, \textit{R\^oles/Affectations}. Si une de ces informations est erron\'ee, cr\'eez un signalement (voir chapitre~\ref{chap:commun}).
\end{warnbox}

\section{Notifications}

La cloche \faBell\ en haut de page indique le nombre de notifications non lues. Les notifications informent l'utilisateur des \'ev\`enements le concernant\,:

\begin{itemize}
  \item R\'eponse \`a une demande de r\^ole (approuv\'ee ou refus\'ee).
  \item Signalement cl\^otur\'e ou trait\'e.
  \item D\'el\'egation reu ou r\'evoqu\'ee.
\end{itemize}

\begin{steps}
  \item Cliquer sur la cloche \faBell\ pour afficher la liste des notifications.
  \item Cliquer sur une notification pour la marquer comme lue.
  \item Naviguer vers la fonctionnalit\'e concern\'ee via le lien dans la notification.
\end{steps}

\screenshot{img_notifications}{Panneau de notifications}{Capturer le panneau de notifications ouvert, montrant plusieurs notifications (certaines lues, certaines non lues). La cloche avec un badge num\'erique doit \^etre visible.}
